\chapter*{出版信息}
\markboth{出版信息}{出版信息}
\addcontentsline{toc}{chapter}{出版信息}

Herbert Amann

Joachim Escher

Birkhäuser Verlag Basel · Boston · Berlin

英译者：Silvio Levy 与 Matthew Cargo（Birkhäuser Verlag，2008）

作者：

Herbert Amann

Institut für Mathematik

Universität Zürich

Winterthurerstr. 190

8057 Zürich

Switzerland

电子邮箱：herbert.amann@math.uzh.ch

Joachim Escher

Institut für Angewandte Mathematik

Universität Hannover

Welfengarten 1

30167 Hannover

Germany

电子邮箱：escher@ifam.uni-hannover.de

德文原版以同一书名由瑞士 Birkhäuser Verlag 出版，© 1999 Birkhäuser Verlag

2000 数学学科分类：26-01, 26A42, 26Bxx, 30-01

美国国会图书馆控制号：2008926303

Die Deutsche Bibliothek 出版书目信息

Die Deutsche Bibliothek 将本出版物列入“Deutsche Nationalbibliografie”；详细的文献数据可在互联网上通过 <http://dnb.ddb.de> 查阅。

ISBN 3-7643-7472-3 Birkhäuser Verlag, Basel – Boston – Berlin

本作品受版权保护。保留一切权利，无论是材料的全部还是部分，特别是翻译、重印、插图使用、朗诵、广播、以缩微胶片或其他方式的复制以及存入数据库的权利。任何形式的使用均须取得版权所有人的许可。

© 2008 Birkhäuser Verlag AG，Basel · Boston · Berlin，P.O. Box 133, CH-4010 Basel, Switzerland

隶属于 Springer Science+Business Media

本书采用由无氯纸浆制成的无酸纸印刷。TCF ∞

德国印刷

ISBN 978-3-7643-7472-3

e-ISBN 978-3-7643-7478-5

987654321

www.birkhauser.ch

\chapter*{译者的话}
\markboth{译者的话}{译者的话}
\addcontentsline{toc}{chapter}{译者的话}

本书是 Herbert Amann 与 Joachim Escher 所著 \textit{Analysis II}（德文原版
\textit{Analysis II}，Birkhäuser Verlag 1999；英译者 Silvio Levy 与 Matthew Cargo，
Birkhäuser Verlag 2008）的中文译本。它与《分析学 I》中文译本配套，合起来覆盖
Amann--Escher 三卷本分析学导论的前两卷。

\subsection*{关于这个译本}

这是一项\textbf{个人自发的翻译工程}，起因很简单：这本书写得好，而中文读者读它
要跨过一道语言的坎。第二卷承接第一卷：在数系、线性代数与一元分析的基础之上，
依次讨论\textbf{一元积分学}（第 VI 章）、\textbf{多元微分学}（第 VII 章）与
\textbf{线积分}（第 VIII 章）——从跳跃连续函数与柯西--黎曼积分出发，经一致收敛、
傅里叶级数与伽马函数，进入连续线性映射、可微性、逆映射与隐函数定理、流形与切空间，
最后落到曲线、普法夫形式、线积分与全纯/亚纯函数。这样一条完整的链条，值得有一份
尽可能贴近原文的中文本来对照阅读。

翻译以\textbf{英译本}为底本，同时参照德文原版的书名与体例。全书按章、节组织，
与原文逐节对应，公式、图表与习题全部保留；章节编号、定理编号（「定理 6.1」）
与公式编号（「(6.1)」）一律沿用原书，以便与英文版、德文版对照查阅。书末的
\textbf{参考文献保留英文原样、未作翻译}，便于按原文献索骥；原书的\textbf{索引
未收入本译本}——索引页码指向原书版次，与中译本重排后的页码无法一一对应。

\subsection*{关于本译本的生成方式（请先读这一段）}

必须把话说明白：\textbf{本译本由 AI（大语言模型）辅助完成}。原文的识别、初译、
排版代码，以及相当一部分核对工作，都是在人机协作下做的——机器出草稿，人定口径、
逐页抽看原书页图、逐轮返工。它\textbf{不是}由某位译者逐字推敲多年的职业译本，也
\textbf{没有经过出版社级别的编校}。

因此请读者务必注意：

\textbf{一、译文可能存在错误。}包括但不限于术语取舍、句子理解、公式与符号的抄录、
编号与交叉引用，以及个别地方与原书不一致的表述。

\textbf{二、不能保证与原著完全一致。}虽已尽量以原书页图逐页核对、把发现的讹误
逐条修掉，但「查过」不等于「查全」，遗漏在所难免。

\textbf{三、请勿把本译本当作学术引用、考试或教学的依据。}凡涉及定义、定理与证明的
准确表述，一律以 Amann--Escher 原著（英译本或德文原版）为准。

一句话：可以拿它当\textbf{对照阅读的拐杖}，但不能拿它当\textbf{标准答案}。

\subsection*{译名与记号}

数学术语以国内通行译法为准，重要的名词首次出现时保留英文原词。全书的术语取舍
与记号约定另有对照表，不在此处赘述。原文中以粗体标示的定义性词（如
\textbf{convex}、\textbf{compact}）在译文中同样以粗体标出。

\subsection*{关于校订}

原书是数字排版，正文与公式的准确字形以\textbf{原书页图}为准：文字层里大量数学
符号（积分号、范数、上下标、分式等）并不可靠，不能直接采信。原书本身的排印瑕疵
一律按原样保留、登记在案，不擅自「改正」；属于识别错误的，按页图修正。因此本译本
中个别地方的写法可能与读者的预期不同，那通常正是原书的样子。

\subsection*{内容问题反馈}

错误既然难免，就希望它被发现、被记下来。译者建了一份在线反馈表「内容问题反馈表」，
用来收集勘误、术语建议与排版问题——发现哪里不对，随手填一下即可：

\begin{center}
\url{https://github.com/JeaESenna/Amann-Analysis-zh/issues}
\end{center}


反馈时若能附上\textbf{章号、页码与原文对照}，处理起来会快得多。谢谢每一位愿意
回头指正的读者。

\subsection*{免责与致谢}

原著版权归 Herbert Amann、Joachim Escher 与 Birkhäuser Verlag 所有，英译版权
归 Silvio Levy 与 Matthew Cargo 所有。本译本仅供个人学习与研究使用，\textbf{不得
用于任何商业用途}。若本书对你有用，请务必去读原著——那是最好的感谢方式。

\vspace{2ex}
\begin{flushright}
Jea神纳\\
2026 年 10 月
\end{flushright}
\chapter*{前言}
\markboth{前言}{前言}
\addcontentsline{toc}{chapter}{前言}

与第一卷一样，第二卷所包含的材料远远多于一门为期一学期的课程所能讲授的内容。这类课程可以略去许多优美而根基扎实的应用，这些应用与数学的许多领域有着广泛的联系。我们当然希望学生能独立地钻研这些材料；教师也许会发现，它们对本科讨论班很有用。

关于本书所呈现材料的概览，请参阅目录与各章的导言。一如从前，我们要强调：做大量的习题对于理解所讲的内容是不可或缺的，而且习题还能补全并扩充正文。

在本卷的写作过程中，我们得益于许多人的帮助。我们特别要感谢我们的朋友兼同事 Pavol Quittner 与 Gieri Simonett。他们不仅一丝不苟地审阅了全部书稿，协助剔除了其中的错误，而且通过他们宝贵的改进建议，对最终版本作出了实质性的贡献。我们还要衷心感谢我们的工作人员，他们仔细通读了全部书稿，并逐一记录了其中的勘误与不准确之处。

我们再次向我们那位“排版完美主义者”致以最衷心的感谢：没有其不倦的努力，这本书绝不会像现在这样美观\footnote{本书正文由 \LaTeX{} 排版，插图由 CorelDRAW! 与 Maple 绘制。}。我们还要感谢 Andreas 帮助解决了硬件与软件方面的问题。

最后，我们感谢 Thomas Hintermann 与 Birkhäuser 出版社，感谢双方良好的合作关系，以及他们对我们的交稿期限的体谅。

苏黎世与卡塞尔，1999 年 3 月

H. Amann、J. Escher

\subsection*{第二版前言}

在这一版中，我们改正了错误，消除了不够精确之处，并简化了若干证明。这些有待改进之处都是由读者提请我们注意的。我们向他们以及我们的同事 H. Crauel、A. Ilchmann 与 G. Prokert 致以衷心的感谢。

苏黎世与汉诺威，2003 年 12 月

H. Amann、J. Escher

\subsection*{英译本前言}

我们很高兴向 Silvio Levy 与 Matt Cargo 致谢，感谢他们将德文原著细致而富有见地地译成英文。他们在翻译过程中高效而愉快的合作，令我们十分感激。

苏黎世与汉诺威，2008 年 3 月

H. Amann、J. Escher
